% EJEMPLO de parcelación POR CLASES CON FECHAS: el piloto de Filosofía 10.° (P4) impreso con fechas de
% prueba (semana 1 = 13 sep 2027; clases lunes y jueves), un receso y un festivo (lunes 18 oct 2027,
% Día de la Raza trasladado por la Ley Emiliani).
% 8 semanas × 2 clases por semana = 16 clases. Sin fechas (reutilizable): para imprimir con fechas,
% \fechainicio{AAAA-MM-DD} (lunes de la semana 1) y \diasclase{lun,jue} en el preámbulo.
% Festivos o eventos: \festivo{motivo} en el lugar de la clase que se pierde.
% Los códigos [D1]… remiten a los desempeños de malla-P4.tex (FIL-10-P4-D1…).
% Compilar: bash curriculo/plantilla/compilar.sh curriculo/filosofia/10/parcelacion-P4.tex
\documentclass[parcelacion]{planeacion}
\materia{filosofia}
\grado{10.°} \periodo{IV} \anio{2027}
\clasesporsemana{2}
\fechainicio{2027-09-13}
\diasclase{lun,jue}
\begin{document}
\begin{parcelacion}
  % ---- Semana 1 · D1 La herencia griega
  \clase{\textbf{Una larga herencia.} Pregunta problematizadora del periodo; Platón y
    Aristóteles.}
    {Lectura de \emph{La escuela de Atenas}: qué señala cada filósofo y por qué. Lluvia de
     preguntas sobre fe y razón.}
    {Reproducción de la obra de Rafael; presentación.}
    {Participación; registro de preguntas en el cuaderno. {\scriptsize[D1]}}
  \clase{\textbf{¿Cómo llegó Aristóteles a la Edad Media?} Neoplatónicos, traducciones del árabe y
    del griego, Toledo, Averroes.}
    {En grupos, línea de tiempo de Platón a Tomás con las distancias entre autores.}
    {Papel para la línea de tiempo; texto guía.}
    {Revisión de la línea de tiempo (coevaluación). {\scriptsize[D1]}}
  % ---- Semana 2 · D2 San Agustín
  \clase{\textbf{San Agustín: la interioridad.} \emph{Confesiones} I, 1; «vuelve a ti mismo».}
    {Lectura guiada de fragmentos e identificación de la tesis de cada uno.}
    {Fragmentos de Agustín con referencia.}
    {Escrito breve: ¿dónde buscamos la verdad? {\scriptsize[D2]}}
  \clase{\textbf{Fe y razón.} «Comprende para creer, cree para comprender».}
    {Conversatorio: ¿se opone la fe a la razón? Posturas, argumentos y contraargumentos.}
    {Preguntas del conversatorio.}
    {Participación en el conversatorio. {\scriptsize[D2]}}
  % ---- Semana 3 · D2 → D3
  \clase[q]{\textbf{Cierre de Agustín.}}
    {Quiz de lectura y retroalimentación inmediata.}
    {Quiz impreso.}
    {Quiz. {\scriptsize[D2]}}
  \clase{\textbf{Tomás de Aquino y la universidad medieval.} La recuperación de Aristóteles.}
    {Exposición dialogada sobre las universidades y la escolástica.}
    {Presentación; mapa de las primeras universidades.}
    {Preguntas de comprensión. {\scriptsize[D3]}}
  % ---- Semana 4 · D3 La quaestio
  \clase{\textbf{La estructura de un artículo de la \emph{Suma teológica}.}}
    {Análisis de un artículo para descubrir sus partes: objeciones, «en sentido contrario»,
     respuesta, solución.}
    {Artículo de la \emph{Suma teológica} (fragmento con referencia).}
    {Esquema de la estructura del artículo. {\scriptsize[D3]}}
  \clase[t]{\textbf{El método de la \emph{quaestio}.}}
    {Taller en parejas: escribir una \emph{quaestio} sobre un problema actual (p.~ej. «¿Es lícito
     usar inteligencia artificial en las tareas?»).}
    {Guía de taller con el modelo de \emph{quaestio}.}
    {Taller evaluativo con criterios (estructura, objeciones, respuesta). {\scriptsize[D3]}}
  \receso{RECESO ESCOLAR DE OCTUBRE}
  \festivo{festivo (Día de la Raza, trasladado al lunes)}
  % ---- Semana 5 · D4 ¿Edad Oscura?
  \clase{\textbf{¿Fue «oscura» la Edad Media?} Origen de la etiqueta.}
    {Preparación del debate: dos grupos reúnen evidencias a favor y en contra.}
    {Fichas de evidencias (universidades, catedrales, traducciones, epidemias).}
    {Fichas de evidencias completas. {\scriptsize[D4]}}
  \clase{\textbf{Debate: ¿Edad Oscura?}}
    {Debate con evidencias y cierre con reflexión escrita.}
    {Fichas de evidencias; rúbrica del debate.}
    {Debate (coevaluación) y reflexión escrita. {\scriptsize[D4]}}
  % ---- Semana 6 · D4 Renacimiento y humanismo
  \clase{\textbf{La imprenta y la difusión de las ideas.}}
    {Lectura de una gráfica sobre la expansión de la imprenta y discusión de sus efectos.}
    {Gráfica de la difusión de la imprenta.}
    {Respuestas a preguntas sobre la gráfica. {\scriptsize[D4]}}
  \clase{\textbf{Humanismo: Erasmo de Rotterdam.} \emph{Elogio de la locura} (1511).}
    {Lectura de un fragmento de Erasmo; cuadro comparativo escolástica--humanismo.}
    {Fragmento de Erasmo; formato de cuadro comparativo.}
    {Cuadro comparativo. {\scriptsize[D4]}}
  % ---- Semana 7 · D5 Utopía
  \clase{\textbf{Tomás Moro y la \emph{Utopía} (1516).} Las ovejas que devoran hombres.}
    {Lectura de fragmentos del libro I y del libro II; qué critica Moro y qué propone.}
    {Fragmentos de la \emph{Utopía}.}
    {Preguntas de interpretación. {\scriptsize[D5]}}
  \clase{\textbf{Proyecto: «Nuestra utopía para el municipio».}}
    {Trabajo en grupos: una medida concreta para mejorar la comunidad, sus razones y sus
     posibles riesgos.}
    {Guía del proyecto.}
    {Avance del proyecto. {\scriptsize[D5]}}
  % ---- Semana 8 · Cierre
  \clase[p]{\textbf{Sustentación de los proyectos.}}
    {Presentación de cada grupo y preguntas del curso.}
    {Proyectos de los grupos.}
    {Sustentación (autoevaluación y coevaluación). {\scriptsize[D5]}}
  \clase[e]{\textbf{Evaluación de periodo.} De la escolástica a la \emph{Utopía}.}
    {Evaluación escrita tipo ICFES (selección y preguntas abiertas).}
    {Evaluación impresa.}
    {Evaluación de periodo. {\scriptsize[D1--D5]}}
\end{parcelacion}
\end{document}
